Vulnerability prioritization methods are routinely validated against CISA's
Known Exploited Vulnerabilities (KEV) catalogue, as if KEV listed every
exploited CVE (publicly catalogued vulnerability). We measure what that assumption costs, using NIST's proposed
Likely Exploited Vulnerabilities (LEV) metric as a case study. LEV compounds a
CVE's history of Exploit Prediction Scoring System (EPSS) scores into a
probability that the CVE has already been exploited; we found no published
empirical evaluation of it. We reimplement LEV, reproduce the white paper's
averages for CVEs published since 2021 to within 7\%, and score it on 105,250
CVEs. Against KEV,
LEV over-predicts by a factor of 17.3 while still ranking CVEs well. Part of
the gap lies in LEV's inputs: LEV multiplies scores across EPSS model releases
that rescale the corpus overnight, which accounts for 27--62\% of its mass,
and EPSS scores for CVEs that no list records rise rather than decay, so LEV
does not saturate as its design assumes. To separate LEV's error from KEV's
blind spots, we audit public exploitation lists for dependence, find that two
of them contain KEV in full, and apply capture-recapture estimation with
explicit rules for when an estimate may be reported. Among 54,612 recent
CVEs, KEV records at most 32\% of the exploited CVEs that three
evidence-citing lists observe. For CVEs
with LEV of at least 0.9, the one stratum in which the estimate is identified,
LEV is about twice as confident as the estimated exploitation rate, rather
than the 4.5 times that KEV alone suggests, and stays within about 2.2 times
when LEV is blinded to KEV listings. A monotone correction restores
calibration a year ahead against the broader ground truth but must be refitted
as CVEs age, and neither LEV nor a learned aggregation of EPSS history
anticipates most new listings. We release code, pinned data where licences allow, and per-CVE
truth bounds so that other prioritization methods can be re-scored against more than
KEV.
